\section{CoM 高さ制御と転倒挙動の関係}
\label{sec:discussion_com_height}

次に，CoM 高さの挙動に着目する．Case4 における結果では，提案手法適用時に DS から SS へ移行する直前および移行直後において，CoM の高さが従来手法と比較して高く保たれていることが確認された．

第 \ref{sec:method_overview} 章で示した LIPM および振り子モデルに基づく理論的考察によれば，重心高さの増大は横方向加速度の低下を通じて，転倒方向への倒れ込みを緩和する効果を持つ．Case4 において観測された CoM 高さの増大は，この理論的予測と整合的であり，SS 移行後における ZMP の安定化に寄与した主要因であると考えられる．

一方，Case3 では沈み込み速度が大きく，CoM 高さ制御が十分に機能する前に転倒挙動が進行してしまった．このことから，提案手法の有効性は，CoM 高さを調整するための時間的余裕が確保できる条件に依存していることが示唆される．
